\section{Numerical correspondence and request cost}
\label{sec:results}
\subsection{Original three-family results}
In the six-session CUDA replication, QB reduces speculative request latency relative to PQ by 3.06\% for Qwen3, 6.17\% for SmolLM2 and 2.83\% for Llama. These are complete-request results for the fixed panel and original cost window. They do not imply that every task benefits or that adding a draft is profitable for every model family.

The separate 72-context replay samples three layer positions, both target phases and query widths one through four. All PQ/QB pairs agree bitwise at QK, softmax and AV. The ordinary multi-query path differs from the repaired paths at QK in the multi-query cases; q=1 cases are common-path controls. These operator results establish correspondence for the captured operands and geometries, not a general theorem about device reduction order.

\begin{figure}[t]
\centering
\begin{tikzpicture}[x=1cm,y=1cm,font=\footnotesize]
\node[anchor=west,font=\bfseries\small] at (0,2.65) {Bitwise-equal contexts / 72};
\foreach \x/\t in {3.30/QK,4.95/Softmax,6.60/AV}\node at (\x,2.15) {\t};
\foreach \y/\t in {1.55/MQ vs PQ,.85/MQ vs QB,.15/PQ vs QB}\node[anchor=west] at (0,\y) {\t};
\foreach \y in {1.55,.85}{\foreach \x in {3.30,4.95,6.60}{
 \draw[draft,line width=.65pt] (\x-.64,\y-.25) rectangle (\x+.64,\y+.25);
 \node[text=draft] at (\x,\y) {18 / 72};}}
\foreach \x in {3.30,4.95,6.60}{\draw[qualified,line width=.65pt,fill=qualified!8] (\x-.64,-.10) rectangle (\x+.64,.40);\node[text=qualified] at (\x,.15) {72 / 72};}
\node[anchor=west,text=aux,font=\scriptsize] at (0,-.55) {18 single-query controls; 54 multi-query contexts};
\end{tikzpicture}

\caption{Fixed-operand numerical correspondence from the original 72-context replay. The graphic distinguishes the repaired PQ/QB pair from ordinary multi-query attention; it does not summarize model-output accuracy.}
\label{fig:numeric}
\end{figure}

\subsection{RTX mapping cost is execution-mode dependent}
Controlled RTX experiments show that actual CUDA Graph capture and replay materially change the relative cost of candidate mappings. For selected geometries, QB's completed-graph ON/OFF ratio is 0.615--0.849. Some QH--QB gaps shrink from roughly 12--13 microseconds to below 2 microseconds after graph replay. This is an execution-mode intervention, not a precise decomposition of CPU launch overhead.

A fixed-logical-length padding intervention also separates a cheap backend path from a valid optimization. Padding logical length 1023 to physical 1024 reaches a three-kernel path but fails the registered reference. The reference-qualified 1023-to-1025 representation remains an eight-kernel path and costs about 1.5\% more. A lower kernel count is therefore neither a numerical qualification nor proof of a physical bottleneck.

In the actual CUDA Graph ON complete Q-representation transaction, compact QH does not establish a practical advantage over the dual-QB control: measured ratios are 1.002216 for q=2 and 1.001598 for q=4. The tested transaction includes copy/transpose, Q upload and completion, graph execution/synchronization, and complete ordinary-output readback; K/V and mask are resident and target buffers preallocated. Natural Q production and KV updates are outside this window. This branch is closed for the tested domain; it is not a proof that all hardware mappings are useless.

\subsection{B570: correct q=2 execution, no practical layout gain}
The B570 operator investigation first found a real boundary. A narrow GQA admission attempt changed the AV evaluation path and failed byte equality at key512/q2, with 7,160 of 8,192 AV values differing despite small absolute error. The candidate stopped before timing. A follow-up preserved the existing noncontiguous AV reduction path using a zero V batch stride and passed all six registered synthetic key256/512 by q1/2/4 coordinates. This establishes the fixed operator domain, not general backend support.

A subsequent real Qwen3 speculative request executes q=2 with GQA32:8 on all 36 target layers. A formal four-input panel then qualifies full endpoints, decision-relevant logits and speculative state traces across R, PQ and QB, and measures 64 requests in four independent processes. The equal-input geometric-mean QB/dual-PQ service ratio is 0.995144, a 0.486\% point reduction. It fails the registered 2\% practical-benefit screen. The QB/target-only service ratio is 2.619763: under the tested policy, drafting costs more than the target-only service it is intended to accelerate.

A one-factor draft CPU-thread parameter intervention (8 to 1) also fails to improve the finite-panel service cost: the DRAFT1/dual-baseline ratio is 1.003994, approximately 0.399\% slower. Source analysis found that the relevant embedding graph requests one task and the graph planner clamps worker count, so the intervention does not establish that fewer workers ran. The negative result does not support a CPU-thread optimization or a physical GPU bottleneck.

\subsection{RTX: speculative economics differ from layout economics}
The new RTX closure experiment uses four previously exposed inputs and one target/draft model pair. The prospective driver and complete-logit references were frozen before qualification and timing. The study preserves historical native endpoints and qualifies all decision-relevant logits and speculative cycle/state traces. Actual caller q=2 and QB physical q=2 are observed across all 36 target layers; PQ uses serial q=1 attention. The formal panel includes 64 measured requests and 64 separately accounted warmups in four independent processes.

The QB/dual-PQ complete-service ratio is 0.989284, corresponding to a 1.072\% point reduction. However, the largest input-level A/A imbalance is 5.208\%, and two process-session aggregates favor PQ. Under the predeclared rule, a practical layout benefit is not established. This is a noise-limited, unresolved small effect—not evidence of equivalence and not a claim that QB is 1.07\% faster.

The same RTX panel gives QB/target-only service ratio 0.582325, or a 41.8\% lower service time for QB than R on these inputs and under this execution policy. This is a separate, useful finite-panel result: the speculative policy repays draft cost on this panel even though the incremental QB-over-PQ effect is unresolved. The target-only arm retains historical ubatch 1, while speculative target verification uses ubatch 4 and the draft uses ubatch 1. It is not an optimized serving ceiling and does not establish a universal speculative-decoding speedup.

\begin{figure}[t]
\centering
\begin{tikzpicture}
\begin{groupplot}[group style={group size=3 by 1,horizontal sep=.70cm},
width=.305\textwidth,height=5.6cm,ymin=.5,ymax=3.5,ytick={1,2,3},
tick label style={font=\footnotesize},title style={font=\small,yshift=12pt},label style={font=\footnotesize,align=center},
axis lines*=left,axis line style={aux,line width=.45pt},tick style={aux},grid=major,grid style={gray!13}]
\nextgroupplot[title={\textbf{A} Local representation},xlabel={$T_{C_{\PQ}}/T_{C_{\QB}}$\\$>1$: QB faster},xmin=.995,xmax=1.09,xtick={1,1.04,1.08},yticklabels={Llama,SmolLM2,Qwen3}]
\addplot[aux,dashed] coordinates {(1,.5)(1,3.5)};
\addplot[only marks,mark=diamond*,mark size=2.7pt,qualified] coordinates {(1.043676000000,3)};
\node[font=\scriptsize,anchor=south] at (axis cs:1.043676000000,3.12){1.044};
\addplot[only marks,mark=diamond*,mark size=2.7pt,qualified] coordinates {(1.068111000000,2)};
\node[font=\scriptsize,anchor=south] at (axis cs:1.068111000000,2.12){1.068};
\addplot[only marks,mark=diamond*,mark size=2.7pt,qualified] coordinates {(1.021848000000,1)};
\node[font=\scriptsize,anchor=south] at (axis cs:1.021848000000,1.12){1.022};
\nextgroupplot[title={\textbf{B} Global serialization},xlabel={$T_{C_{\GS}}/T_{C_{\QB}}$\\$>1$: QB faster},xmin=.95,xmax=2.55,xtick={1,1.5,2,2.5},yticklabels={,,}]
\addplot[aux,dashed] coordinates {(1,.5)(1,3.5)};
\addplot[only marks,mark=diamond*,mark size=2.7pt,qualified] coordinates {(2.143525000000,3)};
\node[font=\scriptsize,anchor=south] at (axis cs:2.143525000000,3.12){2.144};
\addplot[only marks,mark=diamond*,mark size=2.7pt,qualified] coordinates {(1.462530000000,2)};
\node[font=\scriptsize,anchor=south] at (axis cs:1.462530000000,2.12){1.463};
\addplot[only marks,mark=diamond*,mark size=2.7pt,qualified] coordinates {(2.175941000000,1)};
\node[font=\scriptsize,anchor=south] at (axis cs:2.175941000000,1.12){2.176};
\nextgroupplot[title={\textbf{C} Draft value by repair},xlabel={$D_L=T_{M_L}/T_{C_L}$\\$>1$: drafting faster},xmin=.53,xmax=1.27,xtick={.6,.8,1,1.2},yticklabels={,,},legend style={at={(.5,1.01)},anchor=south,draw=none,font=\scriptsize,fill=none},legend columns=3]
\addplot[aux,dashed,forget plot] coordinates {(1,.5)(1,3.5)};
\addplot[only marks,mark=square*,mark size=2.4pt,aux] coordinates {(0.709722260013,3.23) (0.592109590630,2.23) (0.724338602453,1.23)};
\addlegendentry{GS}
\addplot[only marks,mark=o,mark size=2.4pt,target] coordinates {(1.125118952936,3.00) (0.616360145288,2.00) (1.107396133613,1.00)};
\addlegendentry{PQ}
\addplot[only marks,mark=diamond*,mark size=2.4pt,qualified] coordinates {(1.143749000000,2.77) (0.633948000000,1.77) (1.116620000000,0.77)};
\addlegendentry{QB}
\end{groupplot}
\end{tikzpicture}

\caption{Within-panel service ratios. Values from different devices and timing contracts are not pooled. The RTX layout point estimate remains unresolved after accounting for A/A imbalance.}
\label{fig:cost}
\end{figure}
